\section{Discussion and Conclusion}
\label{sec:conclusion}

\subsection{Operational Synthesis and Practical Deployment Implications}
\label{sec:conclusion_implications}

From a transit agency perspective, \texttt{FreqDuet} is designed around three practical deployment constraints that govern real-world bus operations control centers:

\begin{enumerate}
  \item \textbf{Compatibility with standard GTFS, AVL, and stop-level APC feeds:} Unlike corridor mathematical programs requiring real-time origin--destination matrices or stop-platform camera queue sensors, \texttt{FreqDuet} operates exclusively on standard GTFS schedules, vehicle GPS/AVL timestamps, and post-boarding APC counts. Our headway-normalized APC rate transformation (Eq.~\eqref{eq:apc_headway_norm}) resolves a subtle but pervasive instrumentation pitfall: because APC records arrive intermittently at vehicle headways rather than uniform clock ticks, unnormalized minute-binned log-rate filters suffer severe downward Jensen bias. Correcting this arrival-interval scaling before Harmonic-RLS filtering allows agencies to extract clean $20\text{-min}$ diurnal trends and localized positive surge residuals directly from existing AVL/APC infrastructure.

  \item \textbf{Protecting terminal turnaround slack on asymmetric commuting corridors:} Empirical evaluation on Boston MBTA Route~111 illustrates why single-timescale holding and unbudgeted hierarchical RL struggle on radial commuter corridors. During the $07{:}00$ inbound peak ($5.3:1$ directional asymmetry) and $16{:}00$ outbound reversal ($3.6:1$ asymmetry), vehicles running counter-peak serve primarily as mobile fleet supply feeding the peak-direction terminal. Under load-blind holding or unbudgeted RL, holding lightly loaded counter-peak buses to equalize backhaul gaps consumes terminal layover slack and starves peak-direction launches. By combining FIFO ETA-to-launch matching with the backhaul express guard (Eq.~\eqref{eq:budget_proposal}) and differentiable hard action masking (Eq.~\eqref{eq:hard_action_mask}), the Turnaround Budget Ledger guarantees that stop-level holding never cannibalizes peak-direction dispatch readiness—delivering monotonic policy improvement across $200$ training days and preventing the congestion-induced collapse suffered by myopic rolling MPC.

  \item \textbf{Matching signal bandwidth to physical control authority:} Our matched-dimension five-way causal ablation demonstrates that hierarchical temporal abstraction alone (\texttt{NoFreq-HRL} or \texttt{Abl-AllBand}) is insufficient when real-world stop-level APC streams contain substantial Poisson shot noise and localized transfer spikes. Filtering out transient bursts before $15\text{-min}$ quota allocation prevents premature trip-quota exhaustion, while routing positive HF innovations to the stop-event controller targets short $5$--$25\text{ s}$ holds strictly where downstream waiting savings exceed onboard passenger delay (Proposition~\ref{prop:osuna_turnaround}). At the same time, our comparison between heterogeneous \texttt{MBTA-111} and homogeneous high-count \texttt{Tidal-Bimodal} clarifies an actionable boundary condition: when corridor demand is spatially uniform and per-stop arrival counts are high, adaptive forgetting factors $\gamma_f(t)$ or hybrid shock-promotion thresholds can further accelerate upper-level responsiveness to day-wide volume shifts.
\end{enumerate}

\subsection{Limitations and Future Research Directions}
\label{sec:conclusion_limitations}

Several limitations of the present study suggest promising avenues for future research. First, \texttt{FreqDuet} enforces exact daily directional trip quotas with homogeneous fixed vehicle pools and does not yet incorporate dynamic short-turning, deadheading, stop-skipping \citep{rodriguez2023cooperative}, or dynamic standby bus substitution \citep{petit2018substitution}. Integrating discrete short-turn and deadhead actions into the upper slow-clock action space—while extending the FIFO Turnaround Budget Ledger to track multi-terminal circulation—could further unlock capacity on corridors with extreme tidal ratios exceeding $5:1$. Second, link travel times in our formulation are exogenous to holding actions aside from dwell-time coupling; coupling stop-event CMDP holding with transit signal priority at signalized intersections \citep{estrada2016bus} would allow joint optimization of green-wave progression and headway regularity. Third, extending single-corridor Harmonic-RLS decomposition to network-level transfer hubs—where rail arrival waves generate scheduled but jittered multi-route feeder surges—represents a natural multi-corridor generalization.

\subsection{Conclusion}
\label{sec:conclusion_summary}

This paper introduced \texttt{FreqDuet}, a frequency-separated and turnaround-budgeted bi-level constrained reinforcement learning framework for real-time control of non-stationary fixed-pool urban bus corridors. Analytically, we derived the plan-consistent bilateral holding target under time-varying scheduled headways and proved how onboard passenger occupancy and terminal turnaround slack govern the marginal person-seconds trade-off of stop-level holding. Methodologically, \texttt{FreqDuet} pairs a headway-normalized Harmonic-RLS demand decomposer and persistent shock promotion gate with a quota-conserving slow-clock Ensemble SAC planner (stabilized by a Newell-consistent terminal-vanishing quota potential and $6$-step Bellman backups), a causal FIFO Turnaround Budget Ledger, and a hard-masked zero-hold-advantage Constrained Ensemble SAC controller. Across both a high-load bimodal tidal corridor and an empirical Fall~2025 Boston MBTA Route~111 weekday benchmark, \texttt{FreqDuet} reduces net passenger journey time by $14.0$--$14.3\text{ s/pax}$ over unheld operations and $45.3$--$48.8\text{ s/pax}$ over bilateral headway equalization, outperforms single-level DRL and four matched-dimension frequency/budget ablations on real-world APC data, maintains monotonic improvement across $200$ training episodes, and prevents terminal fleet starvation under severe out-of-distribution congestion and demand surges.
